\documentclass[11pt]{article}
\usepackage[margin=1in]{geometry}
\usepackage{amsmath}
\usepackage{booktabs}
\usepackage[T1]{fontenc}
\usepackage{lmodern}

\title{Geant4 Scoring in the MCDC--CubeSat Coupling}
\author{}
\date{30 September 2026}

\begin{document}
\maketitle

\section{Scope and interpretation}

The distribution-mode coupling samples Geant4 source neutrons from the MCDC
surface-current distribution. It writes one Geant4 HDF5 result file per modeled
region (OBC, EPS, ADCS, or communications). Only components configured as scored
sensitive volumes (SVs) contribute to the deposition scores below. All tracks
in a Geant4 event---the source primary and its descendants---contribute when
they make a step inside an SV. The quantities are energy-deposition responses
in the modeled volumes, \emph{not} collected charge or validated bit-upset rates.

Let $e$ denote a Geant4 event, $v$ a scored SV, and $s$ a step in that SV. The
event's physical, unweighted energies are
\begin{align}
 E_{\mathrm{total},ev} &= \sum_{s\in(e,v)} E_{\mathrm{dep},s},\\
 E_{\mathrm{NIEL},ev} &= \sum_{s\in(e,v)} E_{\mathrm{NIEL},s},\\
 E_{\mathrm{ion},ev} &= E_{\mathrm{total},ev}-E_{\mathrm{NIEL},ev}.
\end{align}
Here NIEL is the non-ionizing energy deposit reported by the Geant4 step.
The distinction matters: $E_{\mathrm{ion}}$ is the present SEU-response proxy;
NIEL is retained separately, not treated as generated charge. The scorer
rejects non-finite or materially negative step deposition and NIEL exceeding
the step's total deposition beyond a small tolerance.

\section{Scores present in every run}

The distribution handoff assigns each sampled event a source weight $w_e$.
For each SV, run-integrated energies are weighted sums, for example
\begin{equation}
 S_{\mathrm{ion},v}=\sum_e w_e E_{\mathrm{ion},ev},\qquad
 S_{\mathrm{NIEL},v}=\sum_e w_e E_{\mathrm{NIEL},ev}.
\end{equation}
Weights do not change the physical event energy used for histogram placement
or optional detailed event records. In the present distribution sampler,
$w_e$ is the source-distribution total weight divided by the number of sampled
events. The event-weight interpretation assumes that Geant4 does not split or
roulette descendants into independently weighted branches.

\begin{center}
\begin{tabular}{p{0.44\linewidth}p{0.47\linewidth}}
\toprule
HDF5 dataset & Meaning \\
\midrule
\texttt{component\_edep\_mev} & Weighted total deposited energy per SV (MeV). \\
\texttt{component\_niel\_mev} & Weighted non-ionizing deposition per SV (MeV). \\
\texttt{component\_ionizing\_mev} & Weighted ionizing deposition per SV (MeV); total minus NIEL. \\
\texttt{component\_mass\_kg} & Geant4 mass of each scored SV (kg). \\
\texttt{component\_dose\_gy} & Weighted total deposition divided by SV mass (Gy). \\
\texttt{total\_edep\_mev}, \texttt{dose\_gy} & Existing region-level totals over the scored SVs. \\
\bottomrule
\end{tabular}
\end{center}

The older region-level total-deposition spectrum is also retained:
\texttt{edep\_spectrum\_edges\_mev}, \texttt{edep\_spectrum\_counts},
\texttt{edep\_spectrum\_edep\_mev}, and underflow/overflow counts. It bins
the unweighted event total over the scored SVs, while the deposition stored
in a bin is weighted. It is distinct from the per-SV ionizing-energy spectrum.

\subsection{Depositing-track attribution}

For each SV, the ionizing part of a step is assigned to the \emph{track making
that step}, not necessarily to the incident particle or the particle that
initiated an interaction. The seven columns of
\texttt{component\_species\_ionizing\_mev} follow
\texttt{seu\_species\_names}:
\begin{center}
\begin{tabular}{ll}
\toprule
Column & Track identity \\
\midrule
\texttt{electron\_positron} & Electron or positron \\
\texttt{proton} & Proton (PDG magnitude 2212) \\
\texttt{neutron} & Neutron (PDG magnitude 2112) \\
\texttt{gamma} & Photon \\
\texttt{alpha} & Alpha particle (nuclear PDG code) \\
\texttt{ion\_recoil} & Other nuclear-code particles, including deuterons, tritons, and recoils \\
\texttt{other} & All remaining track types \\
\bottomrule
\end{tabular}
\end{center}
The companion \texttt{component\_primary\_ionizing\_mev} and
\texttt{component\_secondary\_ionizing\_mev} split the same energy by
whether the depositing track has Geant4 parent ID zero or nonzero. Thus the
species columns and the primary/secondary columns are two classifications
of one ionizing-energy total, not additional energy to add to it:
\begin{equation}
 S_{\mathrm{ion},v}=\sum_k S_{\mathrm{ion},v,k}
 =S_{\mathrm{ion},v,\mathrm{primary}}
 +S_{\mathrm{ion},v,\mathrm{secondary}}.
\end{equation}
Species attribution depends in part on production cuts: energy deposited
locally in lieu of producing a low-energy charged secondary is attributed to
the parent track.

\subsection{Per-event, per-SV ionizing-energy distribution}

Each SV has a 150-bin logarithmic histogram of the \emph{unweighted}
$E_{\mathrm{ion},ev}$ from 0.001 to 50 MeV. Its 153 slots are zero,
positive underflow ($0<E_{\mathrm{ion}}<0.001$ MeV), 150 ordinary bins,
and overflow ($E_{\mathrm{ion}}\geq 50$ MeV). Every event contributes one
entry to each SV, including events with zero deposition. The datasets are
\texttt{component\_event\_ionizing\_edges\_mev} and the per-SV arrays
\texttt{component\_event\_ionizing\_count},
\texttt{component\_event\_ionizing\_sumw}, and
\texttt{component\_event\_ionizing\_sumw2}. For a slot $b$,
\begin{equation}
 C_{vb}=\sum_e \mathbf{1}[E_{\mathrm{ion},ev}\in b],\quad
 W_{vb}=\sum_e w_e\mathbf{1}[E_{\mathrm{ion},ev}\in b],\quad
 W^{(2)}_{vb}=\sum_e w_e^2\mathbf{1}[E_{\mathrm{ion},ev}\in b].
\end{equation}
These permit post-processing of weighted exceedance counts for a \emph{chosen}
ionizing-energy threshold; no critical charge or upset decision is built into
the scorer. The stored $\sum w^2$ permits a Geant4-only sampling error for a
threshold aligned with a histogram edge; it is not a full uncertainty estimate
for the coupled calculation.

\subsection{Geant4 sampling uncertainty}

The Geant4 event, not an individual step or secondary track, is the independent
observation. For an SV, let $x_e=w_e E_{\mathrm{ion},ev}$, let
$S=\sum_e x_e$ be its existing weighted ionizing-energy sum, and let
$Q=\sum_e x_e^2$ be the new
\texttt{component\_ionizing\_sum\_sq\_mev2} score. With
$N=\texttt{events\_run}>1$, the estimated standard error of $S$ conditional
on the fixed MCDC source distribution is
\begin{equation}
 \widehat{\mathrm{SE}}(S)
 =\sqrt{\frac{N}{N-1}\left(Q-\frac{S^2}{N}\right)}.
\end{equation}
The same event-level squared sum is recorded for each of the seven depositing
species in \texttt{component\_species\_ionizing\_sum\_sq\_mev2}. Dataset
\texttt{component\_species\_positive\_events} records how many Geant4
events had positive ionizing deposition from each species in each SV; it
counts events, not tracks or steps. A zero count therefore means no such
deposition was observed, not that it is physically impossible. For the
equal-weight distribution sampler, the one-sided 95\% upper bound on the
occurrence probability after zero observations is
$1-0.05^{1/N}\simeq 3/N$. The per-bin $W_{vb}$ and $W^{(2)}_{vb}$ similarly
give the sampling error of a weighted threshold-exceedance sum. These errors
do not include uncertainty in the MCDC source estimate, Geant4 physics
models, or the device response. They cannot be recovered for species-energy
sums in older HDF5 output without rerunning Geant4.

\section{Optional selected-event diagnostics}

Setting \texttt{record\_seu\_events=True} enables four linked tables in
\texttt{seu\_diagnostics}. An event is selected if its \emph{maximum} SV
ionizing deposition is at least
\texttt{diagnostic\_min\_Eion\_mev} (default 0.001 MeV, or 1 keV).
This threshold filters only detailed output; all events still enter the
aggregate scores. A selected event gets one \texttt{event\_sv} row for
\emph{every} scored SV, including SVs with zero deposition.

\begin{center}
\begin{tabular}{p{0.27\linewidth}p{0.64\linewidth}}
\toprule
Table & Recorded fields and interpretation \\
\midrule
\texttt{event\_sv} & Geant4 run/event IDs, SV ID, source weight, physical total/NIEL/ionizing energies, seven species contributions, and primary/secondary ionizing contributions. \\
\texttt{sv\_entry} & One row per SV boundary entry/re-entry, plus a primary starting inside an SV: event/SV and track/parent IDs, PDG code, entry kinetic energy, position and direction, creator process, track vertex position and kinetic energy, and logical volume at the vertex. A track born inside an SV is not counted as an entry. \\
\texttt{nuclear\_birth} & Individual secondary neutron, proton, alpha, deuteron/triton, and other nuclear-ion/recoil births anywhere in the modeled Geant4 event: IDs, PDG code, creation energy and position, vertex logical volume, and creator process. \\
\texttt{em\_secondary\_summary} & Per-selected-event counts and sums of creation kinetic energy for secondary electrons, positrons, and photons; no individual EM birth rows. \\
\bottomrule
\end{tabular}
\end{center}

The run/event IDs identify events within a Geant4 region run, not individual
MCDC history IDs. The detailed records are written via worker-local streams
and included in HDF5 after the Geant4 session closes. The archived input for
run 21455525 used \texttt{record\_seu\_events=False}; its aggregate scores
exist, but its four diagnostic tables cannot be reconstructed afterward.

\section{Transport configuration and output use}

The CubeSat example currently requests \texttt{QGSP\_BIC\_HP} and an
electronics-region EM production range cut of 0.001 mm. One shared
\texttt{ElectronicsRegion} contains the scored SV hierarchies. The range cut
is set individually for gamma, electron, and positron; the proton range cut
is explicitly zero. The region HDF5 file records these requested cuts,
the converted gamma/electron/positron/proton energy thresholds by region
material, the Geant4 version, and the physics-list name. A range cut is not
a track-termination threshold; its effective energy threshold is material
dependent.

By default, \texttt{plot\_seu.py} creates one region overview showing
ionizing and NIEL totals, depositing-species and primary/secondary fractions,
and the fraction of source weight with at least 1 keV of SV ionizing
deposition. It shows one-standard-error bars on ionizing-energy sums and the
1-keV exceedance fractions, plus an event-count/relative-error panel for each
depositing species. The normalized species and ancestry bars remain descriptive
fractions, not confidence intervals for those ratios. An individual SV event-energy histogram is created only when
\texttt{--sv} is requested; when detailed records exist, that plot also shows
species contributions for up to 20 high-deposition selected events and their
secondary-track contribution. The separate
\texttt{plot\_component\_edep.py} displays the older total-edep score.
No per-species event-energy histogram grid is stored.

Finally, the modeled scored volumes are not calibrated memory-cell collection
volumes. Converting this response to a physical SEU rate still requires
device-sensitive geometry, charge-generation/collection modeling, a justified
critical-charge or critical-energy response, and scientific validation of
production cuts and transport physics.

\end{document}
